\section{Conclusion}

The experiments investigated the propagation and application of ultrasound in different materials using 
an echoscope. The measurements of the sound velocity in polyacryl showed a dependence on frequency, 
indicating dispersive behavior of the material. In PVC, the absorption coefficient increased with 
frequency, confirming the expected stronger attenuation of high-frequency ultrasound.
Ultrasonic runtime measurements on the polyacrylic test block enabled the determination of internal 
structures with good agreement to the caliper measurements. In addition, the resolution measurements 
demonstrated that higher frequencies improve the axial resolution due to shorter wavelengths, while 
simultaneously reducing the penetration depth because of increased attenuation. The FFT analysis also 
allowed the thickness of a polyacrylic plate to be determined successfully.
Finally, the Debye-Sears experiment yielded a sound velocity in water consistent with the literature value. 
Overall, the experiments confirmed ultrasound as a reliable and versatile non-destructive method for 
material characterization and distance measurements.